\documentclass[../general/general]{subfiles}

\begin{document}

\chapter{Лекция №10}
\noindent\textit{А.\,М.~Белемук \hfill 17 апреля 2021}
\vspace{1em}


% ===== СЕКЦИЯ 1: Операторы рождения и уничтожения бозонов =====
\section{Вторичное квантование для бозонов}

Сегодня мы начинаем изучение вторичного квантования для системы бозонов. Рассматриваем $N$ нерелятивистских невзаимодействующих бозонных частиц в ящике объёма $V$.

\subsection{Одночастичные состояния}

Одночастичные орбитали --- плоские волны с периодическими граничными условиями:
\begin{equation}
    \varphi_\mathbf{p}(\mathbf{r}) = \frac{1}{\sqrt{V}}\,e^{i\mathbf{p}\cdot\mathbf{r}/\hbar}, \qquad \varepsilon_\mathbf{p} = \frac{\mathbf{p}^2}{2m}
\end{equation}

Импульс $\mathbf{p}$ принимает дискретные значения $\mathbf{p} = \frac{2\pi\hbar}{L}\mathbf{n}$, $\mathbf{n} \in \mathbb{Z}^3$.

\subsection{Операторы рождения и уничтожения}

Вводим для каждого состояния $\mathbf{p}$ пару операторов $\hat{a}_\mathbf{p}$ (уничтожения) и $\hat{a}^\dagger_\mathbf{p}$ (рождения), подчиняющихся \textbf{бозонским коммутационным соотношениям}:
\begin{equation}
    [\hat{a}_\mathbf{p},\,\hat{a}^\dagger_{\mathbf{p}'}] = \delta_{\mathbf{p}\mathbf{p}'}, \qquad
    [\hat{a}_\mathbf{p},\,\hat{a}_{\mathbf{p}'}] = 0, \qquad
    [\hat{a}^\dagger_\mathbf{p},\,\hat{a}^\dagger_{\mathbf{p}'}] = 0
    \label{eq:bose_comm}
\end{equation}

Оператор числа частиц в состоянии $\mathbf{p}$:
\begin{equation}
    \hat{n}_\mathbf{p} = \hat{a}^\dagger_\mathbf{p}\hat{a}_\mathbf{p}, \qquad [\hat{n}_\mathbf{p},\,\hat{a}_\mathbf{p}] = -\hat{a}_\mathbf{p}, \qquad [\hat{n}_\mathbf{p},\,\hat{a}^\dagger_\mathbf{p}] = \hat{a}^\dagger_\mathbf{p}
\end{equation}

\subsection{Вакуумное состояние и построение базисных векторов}

Вакуумное состояние $|0\rangle$ определяется условиями:
\begin{equation}
    \langle 0|0\rangle = 1, \qquad \hat{a}_\mathbf{p}|0\rangle = 0 \quad \forall\,\mathbf{p}
\end{equation}

Это не нулевой вектор, а нормированный вектор, описывающий состояние без частиц. Состояния с $n_\mathbf{p}$ бозонами в орбитали $\mathbf{p}$ строятся последовательным действием оператора рождения:
\begin{equation}
    |\mathbf{p};\,n\rangle = \frac{(\hat{a}^\dagger_\mathbf{p})^n}{\sqrt{n!}}\,|0\rangle
\end{equation}

Общее многочастичное состояние с числами заполнения $\{n_\mathbf{p}\}$ (сколько частиц находится в каждой орбитали):
\begin{equation}
    |\{n_\mathbf{p}\}\rangle = \prod_\mathbf{p} \frac{(\hat{a}^\dagger_\mathbf{p})^{n_\mathbf{p}}}{\sqrt{n_\mathbf{p}!}}\,|0\rangle
    \label{eq:fock_state}
\end{equation}

% ===== КОНЕЦ СЕКЦИИ 1 =====


% ===== СЕКЦИЯ 2: Пространство Фока =====
\section{Пространство Фока}

\subsection{Структура пространства}

Стандартное пространство Гильберта квантовой механики описывает системы с фиксированным числом частиц. Вторичное квантование работает в \textbf{пространстве Фока} --- прямой сумме пространств для любого числа частиц:
\begin{equation}
    \mathcal{F} = \mathcal{B}_0 \oplus \mathcal{B}_1 \oplus \mathcal{B}_2 \oplus \cdots \oplus \mathcal{B}_N \oplus \cdots
\end{equation}

Здесь $\mathcal{B}_N$ --- гильбертово пространство $N$-частичных состояний бозонов. Элементы каждого подпространства:
\begin{itemize}
    \item $\mathcal{B}_0$: вакуум $|0\rangle$
    \item $\mathcal{B}_1$: однoчастичные состояния $\hat{a}^\dagger_\mathbf{p}|0\rangle \equiv |\mathbf{p}\rangle$
    \item $\mathcal{B}_N$: $N$-частичные состояния $|\{n_\mathbf{p}\}\rangle$, $\sum_\mathbf{p} n_\mathbf{p} = N$
\end{itemize}

\subsection{Базис пространства Фока}

Введём сокращённое обозначение: набор всех чисел заполнения $\{n_\mathbf{p}\}$ обозначим $\alpha$. Тогда вектор~\eqref{eq:fock_state} записывается как $|\alpha\rangle$. Эти векторы образуют ортонормированный базис:
\begin{equation}
    \langle\alpha|\beta\rangle = \delta_{\alpha\beta}
\end{equation}

где $\delta_{\alpha\beta} = 1$ тогда и только тогда, когда все числа заполнения совпадают: $n_\mathbf{p}^{(\alpha)} = n_\mathbf{p}^{(\beta)}$ для всех $\mathbf{p}$.

Произвольный вектор состояния $|\Psi\rangle \in \mathcal{F}$ разлагается по этому базису:
\begin{equation}
    |\Psi\rangle = \sum_\alpha c_\alpha\,|\alpha\rangle
\end{equation}

Оператор $\hat{a}_\mathbf{p}$ переводит $\mathcal{B}_N \to \mathcal{B}_{N-1}$, оператор $\hat{a}^\dagger_\mathbf{p}$ --- $\mathcal{B}_N \to \mathcal{B}_{N+1}$.

% ===== КОНЕЦ СЕКЦИИ 2 =====


% ===== СЕКЦИЯ 3: Гамильтониан в пространстве Фока =====
\section{Операторы в пространстве Фока. Одночастичные операторы}

\subsection{Гамильтониан без взаимодействия}

Гамильтониан невзаимодействующей системы:
\begin{equation}
    \hat{H}_0 = \sum_{i=1}^N \hat{h}(\mathbf{r}_i), \qquad \hat{h}(\mathbf{r}) = -\frac{\hbar^2\nabla^2}{2m}
\end{equation}

является \textbf{однoчастичным оператором} (суммой операторов, каждый из которых действует только на одну частицу). Это же справедливо для любого оператора вида $\hat{F} = \sum_i f(\mathbf{r}_i)$.

Утверждение: в представлении вторичного квантования любой однoчастичный оператор $\hat{F}$ записывается как:
\begin{equation}
    \boxed{\hat{F} = \sum_{\mathbf{p},\mathbf{p}'} \langle\mathbf{p}|f|\mathbf{p}'\rangle\,\hat{a}^\dagger_\mathbf{p}\hat{a}_{\mathbf{p}'}}
    \label{eq:one_body}
\end{equation}

где $\langle\mathbf{p}|f|\mathbf{p}'\rangle = \int d^3r\,\varphi^*_\mathbf{p}(\mathbf{r})\,f(\mathbf{r})\,\varphi_{\mathbf{p}'}(\mathbf{r})$ --- матричный элемент однoчастичного оператора.

\subsection{Гамильтониан в собственном базисе}

Если в качестве базиса выбрать собственные функции $\hat{h}$, то матрица $\langle\mathbf{p}|h|\mathbf{p}'\rangle = \varepsilon_\mathbf{p}\delta_{\mathbf{p}\mathbf{p}'}$ диагональна. Формула~\eqref{eq:one_body} даёт:
\begin{equation}
    \hat{H}_0 = \sum_\mathbf{p} \varepsilon_\mathbf{p}\,\hat{a}^\dagger_\mathbf{p}\hat{a}_\mathbf{p} = \sum_\mathbf{p} \varepsilon_\mathbf{p}\,\hat{n}_\mathbf{p}
    \label{eq:H0_diag}
\end{equation}

Это в точности аналог гамильтониана фотонного поля или фононов. Формула~\eqref{eq:H0_diag} действует сразу на всё пространство Фока: не надо думать о фиксированном числе частиц.

Действие на базисный вектор:
\begin{equation}
    \hat{H}_0\,|\{n_\mathbf{p}\}\rangle = \left(\sum_\mathbf{p} \varepsilon_\mathbf{p}\,n_\mathbf{p}\right)|\{n_\mathbf{p}\}\rangle
\end{equation}

\subsection{Пример: оператор полного импульса}

Оператор полного импульса системы $\hat{\mathbf{P}} = \sum_i \hat{\mathbf{p}}_i$ --- также однoчастичный. Матричный элемент:
\begin{equation}
    \langle\mathbf{p}|\hat{\mathbf{p}}_{\text{one}}|\mathbf{p}'\rangle = \int d^3r\,\varphi^*_\mathbf{p}(\mathbf{r})\!\left(-i\hbar\nabla\right)\varphi_{\mathbf{p}'}(\mathbf{r}) = \mathbf{p}'\,\delta_{\mathbf{p}\mathbf{p}'}
\end{equation}

Поэтому:
\begin{equation}
    \hat{\mathbf{P}} = \sum_\mathbf{p} \mathbf{p}\,\hat{a}^\dagger_\mathbf{p}\hat{a}_\mathbf{p}
\end{equation}

% ===== КОНЕЦ СЕКЦИИ 3 =====


% ===== СЕКЦИЯ 4: Полевые операторы =====
\section{Полевые операторы $\hat{\psi}(\mathbf{r})$ и $\hat{\psi}^\dagger(\mathbf{r})$}

\subsection{Определение и смысл}

До сих пор базисом служили плоские волны (собственные функции $\hat{h}$). Перейдём к другому базису: собственным функциям оператора координаты $\hat{\mathbf{r}}$, то есть к состояниям, локализованным в точке $\mathbf{r}$.

Определим \textbf{полевые операторы}:
\begin{equation}
    \hat{\psi}(\mathbf{r}) = \sum_\mathbf{p} \varphi_\mathbf{p}(\mathbf{r})\,\hat{a}_\mathbf{p}, \qquad
    \hat{\psi}^\dagger(\mathbf{r}) = \sum_\mathbf{p} \varphi^*_\mathbf{p}(\mathbf{r})\,\hat{a}^\dagger_\mathbf{p}
    \label{eq:field_ops}
\end{equation}

Это в точности разложение волновой функции по плоским волнам --- но с тем отличием, что коэффициенты $\hat{a}_\mathbf{p}$ заменены на \emph{операторы}. Поэтому и метод называется \emph{вторичным квантованием}: классические коэффициенты разложения «квантуются».

Физический смысл:
\begin{itemize}
    \item $\hat{\psi}^\dagger(\mathbf{r})$ --- оператор рождения частицы в точке $\mathbf{r}$: $\hat{\psi}^\dagger(\mathbf{r})|0\rangle = |\mathbf{r}\rangle$
    \item $\hat{\psi}(\mathbf{r})$ --- оператор уничтожения частицы в точке $\mathbf{r}$
\end{itemize}

\subsection{Коммутационные соотношения полевых операторов}

Из бозонских соотношений~\eqref{eq:bose_comm} и условия полноты $\sum_\mathbf{p}\varphi^*_\mathbf{p}(\mathbf{r})\varphi_\mathbf{p}(\mathbf{r}') = \delta^{(3)}(\mathbf{r}-\mathbf{r}')$ следует:
\begin{equation}
    [\hat{\psi}(\mathbf{r}),\,\hat{\psi}^\dagger(\mathbf{r}')] = \delta^{(3)}(\mathbf{r}-\mathbf{r}')
    \label{eq:field_comm}
\end{equation}
\begin{equation}
    [\hat{\psi}(\mathbf{r}),\,\hat{\psi}(\mathbf{r}')] = 0, \qquad [\hat{\psi}^\dagger(\mathbf{r}),\,\hat{\psi}^\dagger(\mathbf{r}')] = 0
\end{equation}

\textbf{Вывод соотношения~\eqref{eq:field_comm}.} Подставим разложения~\eqref{eq:field_ops}:
\begin{equation}
    [\hat{\psi}(\mathbf{r}),\,\hat{\psi}^\dagger(\mathbf{r}')] = \sum_{\mathbf{p},\mathbf{p}'}\varphi_\mathbf{p}(\mathbf{r})\,\varphi^*_{\mathbf{p}'}(\mathbf{r}')\,[\hat{a}_\mathbf{p},\,\hat{a}^\dagger_{\mathbf{p}'}]
    = \sum_\mathbf{p}\varphi_\mathbf{p}(\mathbf{r})\,\varphi^*_\mathbf{p}(\mathbf{r}') = \delta^{(3)}(\mathbf{r}-\mathbf{r}')
\end{equation}

\subsection{Оператор плотности числа частиц}

Рассмотрим комбинацию $\hat{\psi}^\dagger(\mathbf{r})\hat{\psi}(\mathbf{r})$. Она является \textbf{оператором плотности числа частиц} в точке $\mathbf{r}$:
\begin{equation}
    \hat{\rho}(\mathbf{r}) = \hat{\psi}^\dagger(\mathbf{r})\hat{\psi}(\mathbf{r})
\end{equation}

Оператор полного числа частиц:
\begin{equation}
    \hat{N} = \sum_\mathbf{p}\hat{a}^\dagger_\mathbf{p}\hat{a}_\mathbf{p} = \int d^3r\,\hat{\psi}^\dagger(\mathbf{r})\hat{\psi}(\mathbf{r})
    \label{eq:N_field}
\end{equation}

Коммутационные соотношения с $\hat{N}$:
\begin{equation}
    [\hat{\psi}(\mathbf{r}),\,\hat{N}] = \hat{\psi}(\mathbf{r}), \qquad [\hat{\psi}^\dagger(\mathbf{r}),\,\hat{N}] = -\hat{\psi}^\dagger(\mathbf{r})
\end{equation}

Следствие: $\hat{\psi}(\mathbf{r})$ переводит $N$-частичные состояния в $(N-1)$-частичные, а $\hat{\psi}^\dagger(\mathbf{r})$ --- в $(N+1)$-частичные.

% ===== КОНЕЦ СЕКЦИИ 4 =====


% ===== СЕКЦИЯ 5: Гамильтониан через полевые операторы =====
\section{Гамильтониан через полевые операторы}

\subsection{Однoчастичный гамильтониан}

Подставив разложения~\eqref{eq:field_ops} в выражение для $\hat{H}_0$ и используя~\eqref{eq:one_body}:
\begin{equation}
    \hat{H}_0 = \int d^3r\,\hat{\psi}^\dagger(\mathbf{r})\!\left(-\frac{\hbar^2\nabla^2}{2m}\right)\hat{\psi}(\mathbf{r})
    \label{eq:H0_field}
\end{equation}

Здесь дифференциальный оператор $-\hbar^2\nabla^2/(2m)$ действует на $\hat{\psi}(\mathbf{r})$ как на функцию координаты, а весь оператор $\hat{H}_0$ действует в пространстве Фока.

\textbf{Проверка.} Подействуем~\eqref{eq:H0_field} на состояние $|1_\mathbf{p}\rangle = \hat{a}^\dagger_\mathbf{p}|0\rangle$:

\begin{align}
    \hat{H}_0\,|1_\mathbf{p}\rangle &= \int d^3r\,\hat{\psi}^\dagger(\mathbf{r})\!\left(-\frac{\hbar^2\nabla^2}{2m}\right)\hat{\psi}(\mathbf{r})\,\hat{a}^\dagger_\mathbf{p}|0\rangle \notag\\
    &= \int d^3r\,\hat{\psi}^\dagger(\mathbf{r})\!\left(-\frac{\hbar^2\nabla^2}{2m}\right)\!\left(\varphi_\mathbf{p}(\mathbf{r})\,\hat{a}_\mathbf{p}\hat{a}^\dagger_\mathbf{p}|0\rangle + \hat{a}^\dagger_\mathbf{p}\hat{\psi}(\mathbf{r})|0\rangle\right)
\end{align}

Второе слагаемое обращается в нуль (оператор уничтожения на вакуум). Используя $[\hat{a}_\mathbf{p},\hat{a}^\dagger_\mathbf{p}] = 1$:
\begin{equation}
    \hat{H}_0\,|1_\mathbf{p}\rangle = \int d^3r\,\hat{\psi}^\dagger(\mathbf{r})\!\left(-\frac{\hbar^2\nabla^2}{2m}\right)\varphi_\mathbf{p}(\mathbf{r})\cdot|0\rangle = \varepsilon_\mathbf{p}\int d^3r\,\hat{\psi}^\dagger(\mathbf{r})\,\varphi_\mathbf{p}(\mathbf{r})\,|0\rangle = \varepsilon_\mathbf{p}\,|1_\mathbf{p}\rangle
\end{equation}

Результат совпадает с ожидаемым.

\subsection{Двухчастичный оператор взаимодействия}

Оператор попарного взаимодействия:
\begin{equation}
    \hat{V} = \frac{1}{2}\sum_{i\neq j} v(\mathbf{r}_i - \mathbf{r}_j) = \frac{1}{2}\sum_{i\neq j} f(\mathbf{r}_i,\,\mathbf{r}_j)
\end{equation}

является \textbf{двухчастичным оператором}: он действует одновременно на координаты двух частиц. В представлении вторичного квантования двухчастичные операторы записываются в виде (вывод --- на следующей лекции):
\begin{equation}
    \hat{V} = \frac{1}{2}\int d^3r\,d^3r'\,\hat{\psi}^\dagger(\mathbf{r})\hat{\psi}^\dagger(\mathbf{r}')\,v(\mathbf{r}-\mathbf{r}')\,\hat{\psi}(\mathbf{r}')\hat{\psi}(\mathbf{r})
    \label{eq:V_field}
\end{equation}

Полный гамильтониан взаимодействующего бозонного газа:
\begin{equation}
    \hat{H} = \hat{H}_0 + \hat{V} = \int d^3r\,\hat{\psi}^\dagger(\mathbf{r})\!\left(-\frac{\hbar^2\nabla^2}{2m}\right)\hat{\psi}(\mathbf{r}) + \frac{1}{2}\int d^3r\,d^3r'\,\hat{\psi}^\dagger(\mathbf{r})\hat{\psi}^\dagger(\mathbf{r}')\,v(\mathbf{r}-\mathbf{r}')\,\hat{\psi}(\mathbf{r}')\hat{\psi}(\mathbf{r})
    \label{eq:H_full}
\end{equation}

На следующей лекции мы выведем~\eqref{eq:V_field} и применим полный гамильтониан~\eqref{eq:H_full} к изучению слабо-неидеального бозонного газа.

% ===== КОНЕЦ СЕКЦИИ 5 =====


\end{document}
